%-------------------------------------------------------------------------------
%   CHAPTER 4 - ХЭРЭГЖҮҮЛЭЛТ
%-------------------------------------------------------------------------------

\section{Системийн хэрэгжүүлэлт}
\label{sec:implementation}

Энэ бүлэгт Egüü системийн бодит хэрэгжүүлэлтийг дэлгэрэнгүй танилцуулна. Ашигласан технологиуд, кодын бүтэц, үндсэн функцуудын хэрэгжүүлэлт, мөн туршилтын үр дүнг авч үзнэ.

\subsection{Хөгжүүлэлтийн орчин тохируулах}
\label{subsec:development-setup}

\subsubsection{Шаардлагатай програм хангамжууд}

Egüü төслийг хөгжүүлэхэд дараах програм хангамжууд шаардлагатай:

\begin{table}[H]
\centering
\caption{Хөгжүүлэлтийн орчны шаардлага}
\label{tab:dev-requirements}
\begin{tabular}{|l|l|l|}
\hline
\textbf{Програм} & \textbf{Хувилбар} & \textbf{Зориулалт} \\
\hline
Node.js & 18.x эсвэл дээш & JavaScript runtime \\
\hline
npm/yarn & Latest & Package manager \\
\hline
Expo CLI & Latest & React Native хөгжүүлэлт \\
\hline
VS Code & Latest & Code editor \\
\hline
Git & Latest & Version control \\
\hline
Android Studio & Latest & Android emulator \\
\hline
Xcode & Latest (macOS) & iOS simulator \\
\hline
\end{tabular}
\end{table}

\subsubsection{Төсөл эхлүүлэх}

Expo-аар React Native төсөл үүсгэж, dependencies суулгахын дараа хөгжүүлэлтийн сервер эхлүүлдэг. Нарийвчлалыг Appendix A.8-д үзүүлэв.

\subsection{Firebase тохиргоо}
\label{subsec:firebase-setup}

\subsubsection{Firebase төсөл үүсгэх}

\begin{enumerate}
    \item Firebase Console (https://console.firebase.google.com) руу нэвтрэх
    \item "Add project" товч дарж шинэ төсөл үүсгэх
    \item Төслийн нэр: "eguu-mongolian-script"
    \item Google Analytics идэвхжүүлэх (optional)
    \item iOS болон Android апп бүртгэх
\end{enumerate}

\subsubsection{Firebase SDK суулгах}

Firebase SDK-д багтсан Auth, Firestore, Storage үйлчилгээний packages суулгадаг (Appendix A.8 харна.).

\subsection{Системийн цөм модулиудын програм хангамжийн хэрэгжүүлэлт}
\label{subsec:core-features}

Egüü системийн бүтэц, архитектурын загварын дагуу хөгжүүлэгдсэн цөм модулиудын эх код (Source code) болон хэрэгжүүлэлтийн дэлгэрэнгүй баримтжуулалтыг нээлттэй эхийн GitHub платформ дээр байршуулсан болно: \url{https://github.com/Tugu69ur/movieApp}. Дараах дэд хэсгүүдэд системийн амин чухал функцуудын кодчилол болон техникийн шийдлүүдийг тайлбарлах болно.

\subsubsection{Үүлэн дэд бүтэц буюу Firebase интеграци}

Системийн Backend дэд бүтцийн хувьд серверийн удирдлага шаарддаггүй (Serverless) Firebase платформыг бүрэн хэмжээгээр интеграци хийсэн. Үүнд хэрэглэгчийн бүртгэл, баталгаажуулалтын Authentication, NoSQL өгөгдлийн сан болох Firestore, болон мультимедиа файл хадгалах Storage зэрэг үндсэн үйлчилгээнүүдийг React Native клиент аппликейшнтай холбож, найдвартай бөгөөд хурдан өгөгдөл солилцох сувгийг тохируулсан болно.

\figFirebaseAuth
\figFirebaseDB
\figFirebaseStorage

\subsubsection{Authentication (Баталгаажуулалт)}

\figLoginScreen
\figSignUpScreen

\subsubsection{Spaced Repetition систем}

SM-2 алгоритмыг хэрэгжүүлж, flashcard-ыг хэр сайн санаж байгаагаас хамааран дараагийн давталтын хугацааг тооцоолдог (Appendix A.8 харна.). GitHub дээр дэлгэрэнгүй: \url{https://github.com/Tugu69ur/movieApp/blob/main/utils/spacedRepetition.ts}

\subsubsection{Streak tracking}

Дараалсан өдрүүдийн тоог тооцоолж, мотивацийг үнэлэх функцийг хэрэгжүүлсэн (Appendix A.8 харна.).

\subsection{UI компонентууд}
\label{subsec:ui-components}

React Native болон NativeWind (TailwindCSS) ашиглан орчин үеийн, responsive UI компонентууд бүтээсэн. Үндсэн компонентууд:

\begin{itemize}
    \item \textbf{FlashcardItem:} Flashcard-ын жагсаалтын элемент (зураг, текст, дуртай товч)
    \item \textbf{StreakCounter:} Дараалсан өдрүүдийг харуулах (flame icon, тоо)
    \item \textbf{CategoryPill:} Категори сонгох товчнууд
    \item \textbf{DailyWord:} Өдрийн үгийн карт
\end{itemize}

Бүх компонентын эх код: \url{https://github.com/Tugu69ur/movieApp/tree/main/components}

% UI Component Screenshots
\figHomeScreen
\figSearchScreen
\figProfileScreen
\figSettingsScreen
\figNotificationScreen
\figSecurityPrivacyScreen











\subsection{OCR систем хэрэгжүүлэлт}
\label{subsec:ocr-implementation}

Egüü системийн нэг онцлог функц нь Монгол бичгийг таних OCR (Optical Character Recognition) технологи юм. TensorFlow.js ашиглан хэрэглэгч гараар бичсэн Монгол бичгийг таних боломжтой.

\figCameraInputScreen
\figHandwritingInput
\figHandwrittenText
\figOCRRegularView
\figOCRImageToText
\figOCRCyrillicScript

\subsubsection{OCR загварын гүйцэтгэл}

TensorFlow.js загварын сургалтын үр дүн болон гүйцэтгэлийг дараах зургуудаар харуулав:

\subsubsection{OCR загварын сургалтын дэлгэрэнгүй}

\paragraph{Dataset бүрдүүлэлт:}

Монгол бичгийн OCR системийг сургахын тулд өргөн хүрээтэй dataset бэлтгэсэн. Dataset нь Монгол бичгийн дуу, үг, өгүүлбэрүүдийн зургуудаас бүрдэнэ.

\begin{table}[H]
\centering
\caption{OCR сургалтын dataset-ийн бүрдэл}
\label{tab:ocr-dataset}
\begin{tabular}{|l|r|l|}
\hline
\textbf{Хэсэг} & \textbf{Зургийн тоо} & \textbf{Тайлбар} \\
\hline
Training set & 40,000 & Загварыг сургахад ашигласан \\
\hline
Validation set & 9,615 & Сургалтын явцад үнэлгээ хийхэд ашигласан \\
\hline
Test set & 20,000 & Эцсийн нарийвчлалыг хэмжихэд ашигласан \\
\hline
\textbf{Нийт} & \textbf{69,615} & \\
\hline
\end{tabular}
\end{table}

\paragraph{Монгол бичгийн тэмдэгтүүд:}

Системд дараах 33 үндсэн Монгол бичгийн тэмдэгтийг таних боломжтой:

\begin{itemize}
    \item \textbf{Эгшиг үсгүүд (7):} a, e, i, o, u, ö, ü (уламжлалт Монгол бичгийн тэмдэгтүүд)
    \item \textbf{Гийгүүлэгч үсгүүд (26):} n, ng, b, p, kh, g, m, l, s, sh, t, d, ch, j, y, r, w, f, k, q, ts, dz, h зэрэг уламжлалт бичгийн гийгүүлэгчүүд
    \item \textbf{Хоосон зай:} Үг хоорондын зай
\end{itemize}

\paragraph{CRNN загварын архитектур:}

Монгол бичгийг таних системд CRNN (Convolutional Recurrent Neural Network) архитектурыг ашигласан.

\begin{table}[H]
\centering
\caption{CRNN загварын бүтэц}
\label{tab:crnn-architecture}
\begin{tabular}{|l|l|l|}
\hline
\textbf{Давхарга} & \textbf{Параметр} & \textbf{Зориулалт} \\
\hline
\multicolumn{3}{|c|}{\textbf{CNN хэсэг}} \\
\hline
Conv2D layers (6) & 64→128→256→512 filters & Feature extraction \\
\hline
\multicolumn{3}{|c|}{\textbf{RNN хэсэг}} \\
\hline
Bidirectional LSTM & 256 hidden units & Sequence modeling \\
\hline
\multicolumn{3}{|c|}{\textbf{Output}} \\
\hline
Fully Connected + CTC & 34 classes & Character recognition \\
\hline
\end{tabular}
\end{table}

\paragraph{Сургалтын параметрүүд:}

\begin{itemize}
    \item \textbf{Optimizer:} Adam (learning rate: 1e-4)
    \item \textbf{Loss function:} CTC (Connectionist Temporal Classification)
    \item \textbf{Batch size:} 8
    \item \textbf{Epochs:} 20
    \item \textbf{Image size:} 32×128 pixels
    \item \textbf{Device:} CUDA GPU (эсвэл CPU)
    \item \textbf{Gradient clipping:} 1.0 (тогтвортой байдлыг хангах)
\end{itemize}

\paragraph{Сургалтын үр дүн:}

20 epoch-ийн дараа дараах үр дүнд хүрсэн:

\begin{itemize}
    \item \textbf{Training accuracy:} ~92\%
    \item \textbf{Validation accuracy:} ~87\%
    \item \textbf{Test accuracy:} ~85\%
    \item \textbf{Character Error Rate (CER):} ~15\%
    \item \textbf{Word Error Rate (WER):} ~28\%
\end{itemize}

Загвар нь эхний 5-7 epoch-д хурдан сурч, дараа нь аажмаар сайжирсан. Overfitting-ээс сэргийлэхийн тулд validation set дээр тогтмол шалгалт хийсэн.

\subsection{Туршилт ба тестлэлт}
\label{subsec:testing}

Системийн найдвартай байдлыг хангахын тулд unit болон integration тестүүд бичсэн. Jest болон React Native Testing Library ашигласан.

\paragraph{Тестийн хамрах хүрээ:}

\begin{itemize}
    \item \textbf{Unit тестүүд:} Spaced Repetition алгоритм, streak tracking функцууд
    \item \textbf{Integration тестүүд:} Firebase Firestore үйлдлүүд, Authentication
    \item \textbf{Component тестүүд:} UI компонентуудын render, interaction
\end{itemize}

Бүх тестийн эх код: \url{https://github.com/Tugu69ur/movieApp/tree/main/__tests__}

\subsection{Гүйцэтгэлийн оновчлол}
\label{subsec:performance-optimization}

Апп-ийн гүйцэтгэлийг сайжруулахын тулд дараах оновчлолуудыг хийсэн:

\begin{itemize}
    \item \textbf{Зураг кэшлэх:} Expo Image ашиглан memory-disk кэшлэх
    \item \textbf{Lazy loading:} React.lazy болон Suspense ашиглан компонент ачаалах
    \item \textbf{Firestore offline persistence:} Интернэтгүй үед ч ажиллах
    \item \textbf{Мемо хийх:} React.memo, useMemo, useCallback ашиглан дахин render багасгах
\end{itemize}

Дэлгэрэнгүй: \url{https://github.com/Tugu69ur/movieApp}

\subsection{Бүлгийн дүгнэлт}

Энэхүү бүлэгт Egüü системийн дизайн, архитектурын дагуу хийгдсэн практик кодчилол болон бодит хэрэгжүүлэлтийн үйл явцыг дэлгэрэнгүй тайлбарлалаа. Firebase Authentication, Firestore өгөгдлийн сан, Spaced Repetition танин мэдэхүйн алгоритм, хэрэглэгчийн мотивацийн Streak tracking болон орчин үеийн UI компонентууд зэрэг архитектурын үндсэн элементүүдийн програмын кодыг танилцуулан задлан шинжиллээ. Мөн програмын найдвартай ажиллагааг хангах хэсэгчилсэн (Unit) болон нэгтгэсэн (Integration) тестүүд, аппликейшны гүйцэтгэлийг нэмэгдүүлэх оновчлолын аргуудыг авч үзсэн болно.

Системийн хэрэгжүүлэлтээс гарсан голлох үр дүнгүүд:
\begin{itemize}
    \item React Native болон Expo хүрээг ашиглан iOS болон Android төхөөрөмжүүд дээр тогтвортой ажиллах cross-platform аппликейшныг амжилттай хөгжүүлэн гаргав.
    \item Firebase үүлэн платформ дээр төвлөрсөн, серверийн засвар үйлчилгээ шаардахгүй найдвартай, өргөжүүлэх боломжтой Backend дэд бүтцийг нэвтрүүлсэн.
    \item Танин мэдэхүйд суурилсан SuperMemo-2 spaced repetition (зайтай давталт) алгоритмыг бүрэн кодчилж, тасралтгүй суралцах үйл явцыг программ хангамжаар дэмжиж чадсан.
    \item Хэрэглэгч төвтэй, харилцан үйлчлэл өндөртэй орчин үеийн UI компонентуудыг зохион бүтээж, програмын интерфэйсийг хялбар бөгөөд ойлгомжтой болгосон.
    \item Програм хангамжийн инженерингийн арга зүйн дагуу тестийн процессуудыг автоматжуулан нэвтрүүлж, кодын чанар, серверийн найдвартай ажиллагааг баталгаажуулав.
\end{itemize}

\vfill